\section*{Capítulo \ref{ch:b2:genome-diversification} --- Mutações e diversificação dos genomas}

\begin{solution}{exo:b2:genome-diversification:1}
GAG: silenciosa (continua glutamato), uma transição (A$\to$G). GTA: de
sentido trocado (valina), uma transversão (A$\to$T). TAA: sem sentido
(parada), uma transversão (G$\to$T). Inserção de uma base: um
\omterm{def:b2:genome-diversification:mutation}{deslocamento do quadro de leitura}, que reescreve todos os códons
seguintes.
\end{solution}

\begin{solution}{exo:b2:genome-diversification:2}
$U = 5\times 10^{-10}\times 4.6\times 10^{6} = 2.3\times 10^{-3}$: uma
divisão em cerca de 430 produz uma \omterm{def:b2:genome-diversification:mutation}{mutação}. Em $10^{9}$ células:
$2.3\times 10^{6}$ \omterm{def:b2:genome-diversification:mutation}{mutações} novas por geração.
\end{solution}

\begin{solution}{exo:b2:genome-diversification:3}
\omterm{def:b2:genome-diversification:hgt}{Transformação}: captação de DNA livre de células lisadas; exige uma
receptora competente; transfere fragmentos cromossômicos (genes da
cápsula nos pneumococos). Conjugação: contato célula a célula através
de um pilus; exige um \omterm{def:b2:unicellular-diversity:prokaryote}{plasmídeo} conjugativo na doadora; transfere
\omterm{def:b2:unicellular-diversity:prokaryote}{plasmídeos} (genes de resistência) e, a partir de uma Hfr, genes
cromossômicos. \omterm{def:b2:genome-diversification:hgt}{Transdução}: um bacteriófago que empacotou DNA do
hospedeiro; transfere um fragmento qualquer (generalizada) ou os genes
vizinhos ao sítio do prófago (especializada).
\end{solution}

\begin{solution}{exo:b2:genome-diversification:4}
\omterm{def:b2:genome-diversification:polyploidy}{Autopoliploide}: conjuntos a mais de uma mesma espécie, por um gameta
diploide (a batata tetraploide, muitos trevos autotetraploides).
\omterm{def:b2:genome-diversification:polyploidy}{Alopoliploide}: os conjuntos de duas espécies reunidos por hibridização
e depois duplicados (trigo-mole, algodão, tabaco). O híbrido diploide
AB não tem dois cromossomos iguais, de modo que nenhum consegue parear
na meiose I; os cromossomos segregam ao acaso, e os gametas são
desbalanceados e inviáveis.
\end{solution}

\begin{solution}{exo:b2:genome-diversification:5}
$2\times 3.2\times 10^{9}\times 1.2\times 10^{-8} = 77$ \omterm{def:b2:genome-diversification:mutation}{mutações} novas.
Em sequência codificante: $77\times 0.015 = 1.2$; que mudam um
aminoácido: $1.2\times 0.7 = 0.8$ --- a maioria das crianças carrega
cerca de uma substituição nova de aminoácido.
\end{solution}

\begin{solution}{exo:b2:genome-diversification:6}
Eventos: $m(N - N_0) = 2\times 10^{-8}\times 10^{9} = 20$. Células
mutantes: $mN\ln(N/N_0) = 20\times \ln 10^{6} = 20\times 13.8 = 276$.
$p_0 = e^{-20} \approx 2\times 10^{-9}$: todas as culturas têm
mutantes, de modo que contar as culturas sem nenhum não mede nada. É
preciso $m(N -
N_0) \approx 1$, isto é, $N \approx 5\times 10^{7}$, caso em que
$p_0 \approx
0.37$ e algumas dezenas de culturas dão $m$ com precisão de um fator
dois.
\end{solution}

\begin{solution}{exo:b2:genome-diversification:7}
A uracila não pertence ao DNA, de modo que uma uracila-DNA glicosilase
pode remover toda uracila que encontra, sem ambiguidade. A timina é
uma base normal: depois que a 5-metilcitosina se desamina, o
pareamento incorreto G$\cdot$T é detectado, mas a maquinaria de reparo
não tem como saber qual fita está errada e, na metade das vezes,
``corrige'' o G, fixando uma transição C$\to$T. Os sítios CpG
metilados mutam, portanto, cerca de dez vezes mais depressa que os
outros sítios e, ao longo do tempo evolutivo, foram convertidos em TpG
e CpA: os genomas de mamíferos contêm apenas cerca de um quinto dos CpG
que a composição de bases preveria.
\end{solution}

\begin{solution}{exo:b2:genome-diversification:8}
\qty{46}{kb} por minuto. Posições: 368, 782, 1150 e
\qty{1840}{kb}; distâncias de 414, 368 e \qty{690}{kb}.
\end{solution}

\begin{solution}{exo:b2:genome-diversification:9}
A segunda cópia de um cromossomo pareia com a primeira cópia do outro;
um crossing-over entre elas produz uma cromátide com a cópia~1, a
cópia~2 e a cópia~2$'$ (três genes) e outra apenas com a
cópia~1$'$. Uma pessoa com um cromossomo de um só gene vindo de cada
genitor tem dois genes $\alpha$ em vez de quatro ($-\alpha/-\alpha$):
as cadeias $\alpha$ são produzidas à metade da taxa, as hemácias são
pequenas e pálidas, com anemia leve ou ausente --- é o traço
$\alpha$-talassêmico.
\end{solution}

\begin{solution}{exo:b2:genome-diversification:10}
Gametas: $n = 21$ (um conjunto de cada, A, B, D). Híbridos: AB, 14
cromossomos; ABD, 21. Trigo-mole $\times$ trigo emmer: AABBD, 35
cromossomos --- os conjuntos A e B pareiam, o conjunto D não tem par,
e a planta é em grande parte estéril. Cada duplicação deu a cada
cromossomo um homólogo idêntico, de modo que se formam bivalentes e os
gametas são balanceados; o retrocruzamento com um genitor produz
plantas com um conjunto sem par, cujos gametas são desbalanceados, de
modo que o hexaploide fica reprodutivamente isolado de seus ancestrais
--- uma espécie num único passo.
\end{solution}

\begin{solution}{exo:b2:genome-diversification:11}
Crescimento logístico: $R(t) = N/\bigl(1 + (N - 1)e^{-\gamma Nt}\bigr)$.
Metade da população quando $(N - 1)e^{-\gamma Nt} = 1$: $t = \ln(N -
1)/\gamma N = 20.7/0.5 = \qty{41}{h}$. Sem o antibiótico, as células
sem \omterm{def:b2:unicellular-diversity:prokaryote}{plasmídeo} crescem \qty{5}{\%} mais depressa, de modo que a razão
entre portadoras e não portadoras decai como $e^{-st}$, com $s = 0.05\ln 2 =
\qty{0.035}{h^{-1}}$ por hora (\omterm{def:b2:unicellular-diversity:division}{tempo de geração} de \qty{1}{h}): passar
de \qty{99}{\%} a \qty{1}{\%} leva $\ln(99^2)/0.035 = \qty{260}{h}$,
onze dias --- e a transferência contínua pode manter o \omterm{def:b2:unicellular-diversity:prokaryote}{plasmídeo}
indefinidamente numa frequência baixa, pronto para o próximo
tratamento com o fármaco.
\end{solution}

\begin{solution}{exo:b2:genome-diversification:12}
O \omterm{thm:b2:genome-diversification:luria}{teste de flutuação} mostrou que as \omterm{def:b2:genome-diversification:mutation}{mutações} de resistência surgem
antes e na ausência do agente seletivo, em momentos aleatórios, de
modo que a \omterm{def:b2:genome-diversification:mutation}{mutação} é cega ao que seria útil; o plaqueamento em réplica
mostrou o mesmo sem expor as células selecionadas em momento algum. A
evolução não é aleatória porque a seleção tria essas variantes cegas
por suas consequências, de modo consistente e cumulativo. O único
ponto em que a própria aleatoriedade é moldada é a taxa: as linhagens
com erros demais ou de menos perdem, de modo que a \omterm{thm:b2:genome-diversification:rate}{taxa de mutação} é
um compromisso adaptado --- mas o que cada \omterm{def:b2:genome-diversification:mutation}{mutação} faz continua
imprevisível e não escolhido.
\end{solution}

\begin{solution}{pb:b2:genome-diversification:1}
\textbf{1.} Soma 119, média $5.95$; média dos quadrados $11489/20 = 574$;
variância $574 - 35 = 540$.
\textbf{2.} Soma 324, média $16.2$; média dos quadrados $5462/20 = 273$;
variância $273 - 262 = 11$.
\textbf{3.} As amostras de uma só cultura seguem Poisson (variância
próxima da média): são sorteios aleatórios de uma mesma fração fixa de
mutantes. As culturas independentes têm uma variância noventa vezes
maior que a média: seus mutantes surgiram em momentos diferentes e
aleatórios durante o crescimento, e quanto mais precoce a \omterm{def:b2:genome-diversification:mutation}{mutação},
maior seu clone.
\textbf{4.} Um prêmio acumulado: uma \omterm{def:b2:genome-diversification:mutation}{mutação} numa das primeiras
divisões, cujo clone cresceu depois com a cultura até cem células.
\textbf{5.} Quatorze de vinte: $p_0 = 0.70$; $m(N - N_0) = -\ln 0.70
= 0.36$; $m = 0.36/2\times 10^{8} = \qty{1.8e-9}{}$ por divisão.
\textbf{6.} $1.8\times 10^{-9}\times 2\times 10^{8}\times \ln(2\times
10^{6}) = 0.36\times 14.5 = 5.2$ células mutantes; observado $5.95$
--- concordância dentro do ruído de vinte culturas.
\textbf{7.} Nenhuma amostra estava vazia (todas compartilham os
mutantes da cultura de origem), de modo que $p_0 = 0$ não dá nenhuma
estimativa; as amostras medem a fração de mutantes de uma única
cultura, que depende de quando suas \omterm{def:b2:genome-diversification:mutation}{mutações} ocorreram.
\textbf{8.} $u = m/20 = \qty{9e-11}{}$ por par de bases por replicação.
\textbf{9.} $U = 9\times 10^{-11}\times 4.6\times 10^{6} =
\qty{4.1e-4}{}$.
\textbf{10.} $10^{12}$ células $\times$ $4.1\times 10^{-4}$ $= 4.1\times
10^{8}$ \omterm{def:b2:genome-diversification:mutation}{mutações} novas.
\textbf{11.} $3\times 4.6\times 10^{6} = 1.4\times 10^{7}$ mudanças
possíveis; cada uma surge $4.1\times 10^{8}/1.4\times 10^{7} \approx 30$
vezes.
\textbf{12.} Toda substituição de base única, inclusive cada uma das
que conferem resistência a um fármaco vencido por uma única \omterm{def:b2:genome-diversification:mutation}{mutação},
está presente em dezenas de células antes de o fármaco ser adicionado;
o fármaco apenas as seleciona.
\textbf{13.} Duas \omterm{def:b2:genome-diversification:mutation}{mutações} independentes na mesma célula ocorrem a
$m^2 \approx (2\times 10^{-9})^2 = 4\times 10^{-18}$ por divisão: em
$10^{12}$ células, $4\times 10^{-6}$ por geração --- uma célula
duplamente resistente praticamente nunca preexiste, e é por isso que
a tuberculose é tratada com vários fármacos ao mesmo tempo.
\textbf{14.} Com $a = N - 1$ e $E = e^{-\gamma Nt}$: $R = N/(1 +
aE)$, $\mathrm{d}R/\mathrm{d}t = \gamma N\cdot NaE/(1 + aE)^2$; e
$\gamma R(N - R) = \gamma\,\frac{N}{1 + aE}\cdot\frac{NaE}{1 + aE}$,
o mesmo; $R(0) = N/(1 + N - 1) = 1$.
\textbf{15.} $(N - 1)E = 1$: $t = \ln(N - 1)/\gamma N = 20.7/0.5 =
\qty{41}{h}$.
\textbf{16.} $1 + aE = 1/0.99$, $aE = 0.0101$: $t = \ln(10^{9}/0.0101)
/0.5 = (20.7 + 4.6)/0.5 = \qty{51}{h}$.
\textbf{17.} $\gamma (N/2)^2 = \gamma N\cdot N/4 = 0.5\times 2.5\times
10^{8} = 1.25\times 10^{8}$ transferências por hora.
\textbf{18.} Taxas de crescimento $\ln 2 = \qty{0.693}{h^{-1}}$ e
\qty{0.762}{h^{-1}}; a razão mutante/sensível cresce como $e^{st}$
com $s = \qty{0.069}{h^{-1}}$; passar de $1/N$ a 1 leva $\ln N/s =
20.7/0.069 = \qty{300}{h}$, doze dias.
\textbf{19.} A transferência é proporcional ao produto de portadoras e
não portadoras e ocorre ao ritmo dos encontros, não das divisões; a
descendência é limitada pela vantagem de crescimento do mutante, que é
pequena. Um \omterm{def:b2:unicellular-diversity:prokaryote}{plasmídeo} atravessa a população em dois dias; um mutante
favorecido, em duas semanas.
\textbf{20.} $U = 4.1\times 10^{-2}$; nocivas: $0.041\times 0.88\times
0.3 = 0.011$ por geração.
\textbf{21.} Média $2.2$: $e^{-2.2} = 0.11$, um descendente em cada
nove ileso.
\textbf{22.} Tipo selvagem: $1.1\times 10^{-4}$ por geração, $0.022$
em 200: $e^{-0.022} = 0.98$.
\textbf{23.} Média de $520$ células mutantes; $p_0 = e^{-36} \approx
10^{-16}$: nenhuma cultura vazia.
\textbf{24.} Sob seleção intensa e variável (os antibióticos em
rodízio de um hospital), o suprimento cem vezes maior de variantes
novas supera a carga, e as mutadoras pegam carona nas resistências que
geram. Num ambiente estável, não há nada a ganhar e uma \omterm{def:b2:genome-diversification:mutation}{mutação} nociva
a cada cem gerações a perder, de modo que as linhagens mutadoras são
eliminadas.
\textbf{25.} $m = \qty{1.8e-9}{}$ por divisão para a resistência, $u =
\qty{9e-11}{}$ por par de bases por replicação,
$U = \qty{4.1e-4}{}$ \omterm{def:b2:genome-diversification:mutation}{mutações} por genoma por replicação; o \omterm{def:b2:unicellular-diversity:prokaryote}{plasmídeo}
alcança metade do intestino em \qty{41}{h} e \qty{99}{\%} dele em
\qty{51}{h}.
\end{solution}
